\chapter{Machine-Readable Extraction Layer}
\label{ch:extraction-layer}

This chapter documents the extraction process that transforms source
documents into machine-readable formats, providing the foundation for
automated validation and requirements traceability.

% =============================================================================
\section{Extraction Process Overview}
\label{sec:extraction-overview}

Source documents were processed using automated extraction scripts to
produce machine-readable formats suitable for:

\begin{itemize}[noitemsep]
  \item Text analysis and requirement derivation
  \item Schema validation and data quality checks
  \item RAG (Retrieval-Augmented Generation) pipeline ingestion
  \item Version control and change tracking
\end{itemize}

\begin{table}[h]
\centering
\caption{Extraction Scripts and Outputs}
\label{tab:extraction-scripts}
\begin{tabularx}{\textwidth}{@{}L{3.5cm}L{4cm}L{6.5cm}@{}}
\toprule
\textbf{Script} & \textbf{Source Types} & \textbf{Output Format} \\
\midrule
\texttt{extract\_docs.py} & .doc, .docx & Markdown with YAML frontmatter \\
\texttt{extract\_pdf.py} & .pdf (BPMN) & Markdown + BPMN elements JSON \\
\texttt{extract\_xlsx.py} & .xlsm workbooks & YAML schemas + CSV samples \\
\texttt{build\_rag\_chunks.py} & All extracted text & JSONL chunks for RAG \\
\bottomrule
\end{tabularx}
\end{table}

% =============================================================================
\section{Extracted Document Mapping}
\label{sec:extracted-mapping}

\begin{longtable}{@{}C{2cm}L{5cm}L{6.5cm}@{}}
\caption{Source to Extracted Document Mapping}
\label{tab:extracted-mapping} \\
\toprule
\textbf{Source ID} & \textbf{Extracted File} & \textbf{Extraction Notes} \\
\midrule
\endfirsthead
\toprule
\textbf{Source ID} & \textbf{Extracted File} & \textbf{Extraction Notes} \\
\midrule
\endhead
\srcref{TB-BC-001} & \texttt{business-case-trial-balance.md} & Full text extraction; tables converted to markdown; 1,172 words \\
\srcref{TB-BC-002} & \texttt{business-case-bss-risk-assessment.md} & Full text extraction; 8,908 bytes \\
\srcref{TB-DOI-001} & \texttt{doi-trial-balance-process.md} & Text + 4 images extracted to \texttt{doi-images/} \\
\srcref{TB-BPMN-001} & \texttt{bpmn-tb-review-glo.md} + \texttt{bpmn-elements.json} & Process text + structured BPMN elements \\
\srcref{TB-WB-001} & \texttt{hkg-tb-schema.yaml} + \texttt{hkg-tb-sample.csv} & Schema for 13 sheets; 100-row sample \\
\srcref{TB-WB-002} & \texttt{hkg-pl-schema.yaml} + \texttt{hkg-pl-sample.csv} & Schema for 14 sheets; 100-row sample \\
\bottomrule
\end{longtable}

% =============================================================================
\section{Extraction File Inventory}
\label{sec:extraction-inventory}

\subsection{Extracted Text Files}

\begin{longtable}{@{}L{6cm}C{2cm}C{2cm}L{3.5cm}@{}}
\caption{Extracted Text File Inventory}
\label{tab:extracted-text-inventory} \\
\toprule
\textbf{File Path} & \textbf{Size} & \textbf{Format} & \textbf{Source Ref} \\
\midrule
\endfirsthead
\toprule
\textbf{File Path} & \textbf{Size} & \textbf{Format} & \textbf{Source Ref} \\
\midrule
\endhead
\texttt{extracted/business-case-trial-balance.md} & 9,824 B & Markdown & \srcref{TB-BC-001} \\
\texttt{extracted/business-case-bss-risk-assessment.md} & 8,908 B & Markdown & \srcref{TB-BC-002} \\
\texttt{extracted/doi-trial-balance-process.md} & 865 B & Markdown & \srcref{TB-DOI-001} \\
\texttt{extracted/bpmn-tb-review-glo.md} & 9,154 B & Markdown & \srcref{TB-BPMN-001} \\
\texttt{extracted/bpmn-elements.json} & 2,264 B & JSON & \srcref{TB-BPMN-001} \\
\bottomrule
\end{longtable}

\subsection{Extracted Image Files}

\begin{longtable}{@{}L{6cm}C{2cm}C{2cm}L{3.5cm}@{}}
\caption{Extracted Image File Inventory}
\label{tab:extracted-image-inventory} \\
\toprule
\textbf{File Path} & \textbf{Size} & \textbf{Format} & \textbf{Source Ref} \\
\midrule
\endfirsthead
\toprule
\textbf{File Path} & \textbf{Size} & \textbf{Format} & \textbf{Source Ref} \\
\midrule
\endhead
\texttt{extracted/doi-images/doi-image-001.emf} & 8,148 B & EMF & \srcref{TB-DOI-001} \\
\texttt{extracted/doi-images/doi-image-002.emf} & 8,272 B & EMF & \srcref{TB-DOI-001} \\
\texttt{extracted/doi-images/doi-image-003.png} & 43,650 B & PNG & \srcref{TB-DOI-001} \\
\texttt{extracted/doi-images/doi-image-004.emf} & 8,132 B & EMF & \srcref{TB-DOI-001} \\
\bottomrule
\end{longtable}

\subsection{Training Data Files}

\begin{longtable}{@{}L{6cm}C{2cm}C{2cm}L{3.5cm}@{}}
\caption{Training Data File Inventory}
\label{tab:training-data-inventory} \\
\toprule
\textbf{File Path} & \textbf{Size} & \textbf{Format} & \textbf{Source Ref} \\
\midrule
\endfirsthead
\toprule
\textbf{File Path} & \textbf{Size} & \textbf{Format} & \textbf{Source Ref} \\
\midrule
\endhead
\texttt{training-data/hkg-tb-schema.yaml} & 20,228 B & YAML & \srcref{TB-WB-001} \\
\texttt{training-data/hkg-tb-sample.csv} & 17,916 B & CSV & \srcref{TB-WB-001} \\
\texttt{training-data/hkg-pl-schema.yaml} & 16,982 B & YAML & \srcref{TB-WB-002} \\
\texttt{training-data/hkg-pl-sample.csv} & 19,110 B & CSV & \srcref{TB-WB-002} \\
\bottomrule
\end{longtable}

% =============================================================================
\section{Extraction Frontmatter Format}
\label{sec:extraction-frontmatter}

Each extracted markdown file includes YAML frontmatter for traceability
and metadata validation:

\begin{lstlisting}[language=yamlx,caption=Extraction Frontmatter Example]
---
source_file: 'Business case - Trial Balance (2).doc'
source_type: 'doc'
extracted_at: '2026-04-13T06:16:45.446975+00:00'
word_count: 1172
source_ref: 'TB-BC-001'
---
# Business Case: Trial Balance

## 1. Executive Summary
...
\end{lstlisting}

\paragraph{Frontmatter Fields:}

\begin{table}[h]
\centering
\caption{Extraction Frontmatter Schema}
\label{tab:frontmatter-schema}
\begin{tabularx}{\textwidth}{@{}L{3cm}C{2cm}L{8.5cm}@{}}
\toprule
\textbf{Field} & \textbf{Type} & \textbf{Description} \\
\midrule
\texttt{source\_file} & string & Original filename \\
\texttt{source\_type} & enum & File extension (.doc, .docx, .pdf, .xlsm) \\
\texttt{extracted\_at} & ISO 8601 & Timestamp of extraction \\
\texttt{word\_count} & integer & Word count of extracted text \\
\texttt{source\_ref} & string & Traceability reference ID \\
\bottomrule
\end{tabularx}
\end{table}

% =============================================================================
\section{BPMN Elements Extraction}
\label{sec:bpmn-extraction}

The BPMN process map (\srcref{TB-BPMN-001}) was parsed into structured
JSON to enable automated workflow validation:

\begin{lstlisting}[language=jsonx,caption=BPMN Elements Structure]
{
  "workflowId": "tb-review-glo",
  "version": "1.0.0",
  "states": [
    {
      "stateId": "START",
      "type": "start_event",
      "name": "Prepare Schedule Work Day",
      "transitions": [{"target": "ROLL_FORWARD"}]
    },
    {
      "stateId": "ROLL_FORWARD",
      "type": "task",
      "name": "Roll Forward Prior Month File",
      "timing": "M-3",
      "systems": ["MS Excel", "Shared Drive"],
      "transitions": [{"target": "DOWNLOAD_PSGL"}]
    }
  ],
  "gateways": [...],
  "dataObjects": [...],
  "participants": [...]
}
\end{lstlisting}

\paragraph{Extracted BPMN Statistics:}

\begin{table}[h]
\centering
\caption{BPMN Extraction Statistics}
\label{tab:bpmn-statistics}
\begin{tabular}{@{}lr@{}}
\toprule
\textbf{Element Type} & \textbf{Count} \\
\midrule
Start Events & 1 \\
Task States & 21 \\
Gateway States & 5 \\
Error-Handling States & 3 \\
End Events & 1 \\
\midrule
\textbf{Total States} & \textbf{27} \\
\midrule
Data Objects & 6 \\
Participants & 3 \\
\bottomrule
\end{tabular}
\end{table}

% =============================================================================
\section{Workbook Schema Extraction}
\label{sec:workbook-schema}

Excel workbooks were processed to extract structural schemas without
copying sensitive data:

\subsection{HKG TB Review Schema (\srcref{TB-WB-001})}

\begin{table}[h]
\centering
\caption{HKG TB Review Sheet Inventory}
\label{tab:hkg-tb-sheets}
\begin{tabularx}{\textwidth}{@{}L{4cm}C{2cm}L{7.5cm}@{}}
\toprule
\textbf{Sheet Name} & \textbf{Fields} & \textbf{Purpose} \\
\midrule
IFRS Summary BS & -- & IFRS balance sheet summary view \\
Count of Comments & 7 & Commentary tracking (reviewer, status, hours) \\
BS Variance & 5 & Balance sheet variance analysis with thresholds \\
PL Variance & 4 & P\&L variance analysis \\
GCOA & 5 & Global Chart of Accounts definitions \\
Dept Mapping & 12 & Department-to-vertical mapping \\
Base file & 13 & Account-level base data with mappings \\
\bottomrule
\end{tabularx}
\end{table}

% =============================================================================
\section{Extraction Quality Metrics}
\label{sec:extraction-quality}

\begin{table}[h]
\centering
\caption{Extraction Quality Metrics}
\label{tab:extraction-quality}
\begin{tabularx}{\textwidth}{@{}L{5cm}C{2.5cm}L{6cm}@{}}
\toprule
\textbf{Metric} & \textbf{Value} & \textbf{Human Validation Required} \\
\midrule
Source documents processed & 6 of 6 & Verify all documents included \\
Total extracted text size & 522 KB & Verify no truncation occurred \\
Tables extracted & 12 & Verify table structure preserved \\
Images extracted & 4 & Verify images are referenced correctly \\
BPMN elements parsed & 27 states & Verify state machine completeness \\
Workbook sheets processed & 27 & Verify key sheets included \\
\bottomrule
\end{tabularx}
\end{table}

\begin{hitlbox}{H2 -- Extraction Accuracy Verification}
The human reviewer must verify the following for checkpoint H2:
\begin{enumerate}[noitemsep]
  \item Sample extracted text against original documents (minimum 3 paragraphs per document)
  \item Table formatting is preserved in markdown conversion
  \item BPMN state names and transitions match the source PDF
  \item Workbook schema field names match column headers
  \item No content truncation or corruption occurred
\end{enumerate}
\end{hitlbox}